\section{Exact one-switch minimax on an arbitrary grid}\label{sec:finite-one}
The effect of a grid is not an additive function of the maximum cell width.
It depends on the locations of particular grid points, and uniform relaxed
controls need not be worst-case inputs even on a uniform grid.
This section determines the grid minimax value $F^{\mathcal T}_{n,1}$ for
every grid and every $n\ge3$. Throughout, write
$\mathcal T=(t_0,\ldots,t_N)$, where $0=t_0<\cdots<t_N=T$.

\subsection{An explicit formula and compressed extremizers}
Let $c=\min\{j:t_j\ge T/3\}$ and define
\begin{equation}\label{eq:grid-H}
 H=\min\left\{\frac{T+t_c}{4},\frac{T-t_{c-1}}2\right\}.
\end{equation}
For each $1\le a\le b\le N$, abbreviate
$A=t_a$, $B=t_b$, $P=t_{a-1}$, $Q=t_{b-1}$ and put
\begin{align}
 L_{ab}&=\max\left\{\frac T3,\frac{(n-1)T}{n}-B,
                \frac{(n-1)T-A-(n-1)B}{n}\right\},\label{eq:grid-L}\\
 U_{ab}&=\min\left\{A,\frac{(n-2)A+B}{n},\frac{(n-1)T}{n}-P,\right.\notag\\
 &\hspace{23mm}\left.\frac{(n-1)T-P-(n-1)Q}{n},
          \frac{T-P+(n-2)A}{n}\right\}.\label{eq:grid-U}
\end{align}
Call $(a,b)$ admissible when
\begin{equation}\label{eq:grid-admissible}
 (n-2)(T-A)\le(n-1)B,\qquad L_{ab}\le U_{ab}.
\end{equation}

\begin{theorem}[Arbitrary-grid one-switch minimax]\label{thm:finite-one}
For $n\ge3$ and every strictly increasing finite grid,
\begin{equation}\label{eq:grid-minimax-formula}
 F^{\mathcal T}_{n,1}=\max\{H,\ \max_{(a,b)\text{ admissible}}U_{ab}\},
\end{equation}
where an empty inner maximum is ignored. The relaxed maximizer can be chosen
with at most three time phases and two types of component: either one
mode and $n-1$ identical modes, or two identical modes and $n-2$ identical
modes. For rational grid data, the value and this compressed maximizer
require $O(N^2)$ rational arithmetic operations and have polynomial bit
complexity in the grid encoding length and $\log n$.
\end{theorem}
The theorem allows $N=1$, zero-length phases, and constant integer controls;
by Lemma~\ref{lem:endpoint} it also covers arbitrary measurable relaxed
inputs. The compressed representation lists each component type once with
its multiplicity; writing out all $n$ component functions takes time at
least linear in $n$.

\begin{proof}
We first derive a finite family of closed linear programs. Order the terminal
masses $m_1\ge\cdots\ge m_n$. For fixed initial mode $p$ and switching
time $\tau$, a largest-mass mode other than $p$ is an optimal final mode:
in \eqref{eq:one-switch-error}, this choice reduces the final deficit,
leaves the initial deficit unchanged, and cannot increase the largest
omitted mass. Thus only $p\to1$ with $p\ne1$ and $1\to2$ need be
considered.

For $1\le a\le b\le N$, use variables $u_j,v_j,m_j$ ($j=1,2,3$) and
$E$, where component $3$ is repeated $n-2$ times. With weights
$w=(1,1,n-2)$, impose
\begin{subequations}\label{eq:grid-LP}
\begin{align}
 &0\le u_j\le v_j\le m_j\quad(j=1,2,3),\qquad
 m_1\ge m_2\ge m_3,\qquad E\ge T/3,\\
 &\sum_jw_ju_j=t_a,\quad\sum_jw_jv_j=t_b,
       \quad\sum_jw_jm_j=T,\\
 &T-t_a\le E+m_1\le T-t_{a-1},\qquad
 T-t_b\le E+m_2\le T-t_{b-1},\\
 &u_2+E\le t_a,\qquad v_1+E\le t_b.
\end{align}
\end{subequations}
The \emph{all-initial-modes family} also imposes $u_3+E\le t_a$.
The \emph{two-large-modes family} instead imposes $m_2\ge E$.
Both maximize $E$; an infeasible program is ignored. In particular, the
first family does not impose $m_2\le E$.

Every feasible point gives a valid relaxed cumulative control by linear
interpolation through $0,t_a,t_b,T$. Nonnegative increments and the sum
constraints give nonnegative rates summing to one. Equal times force equal
states, so empty phases can be removed. For $p\to1$, a grid switch before
$t_a$ has final deficit
$T-m_1-\tau\ge T-m_1-t_{a-1}\ge E$.
A switch at or after $t_a$ has initial deficit at least $t_a-u_p$, since
$t-A_p(t)$ is nondecreasing. The same argument for $1\to2$ uses $t_b,v_1$.
In the first family this obstructs every schedule by dominance. In the
second family any schedule omitting mode $1$ or $2$ has error at least
$m_2\ge E$, and the two remaining orders are obstructed by the common
constraints. Constant schedules are included. Hence each feasible point
is a lower-bound witness of value $E$.

Conversely, let $F=F^{\mathcal T}_{n,1}$, attained by compactness.
Giving three modes mass $T/3$ proves $F\ge T/3$. Suppose $F>T/3$ and
fix $T/3<E<F$. For $j=1,2$, let $k_j$ be the first grid index satisfying
$t_{k_j}\ge T-E-m_j$. Both indices are positive: $k_1=0$ would give
the constant mode-$1$ schedule error $T-m_1\le E$. Set $a=k_1\le k_2=b$.
There are at most two masses exceeding $E$.
If $m_2\le E$, every $p\to1$ and $1\to2$ omits only masses at most
$E$. At its first eligible grid time its final deficit is at most $E$,
so its initial deficit must exceed $E$.
If $m_2>E$, apply this reasoning only to $2\to1$ and $1\to2$; their
omitted masses are at most $E$ since $3E>T$.
Averaging modes $3,\ldots,n$ preserves all required inequalities and yields
a feasible point in the respective closed program. Let $E\uparrow F$,
choose a fixed family and pair $(a,b)$ along a subsequence, and take a
convergent subsequence of the bounded endpoint data. This gives a point
with value $F$.
If $F=T/3$, take $E=T/3$, $a=b=c$, $m_1=m_2=E$,
$m_3=E/(n-2)$, $u_1=v_1=u_2=v_2=(t_c-E)/2$, and
$u_3=v_3=E/(n-2)$. We have $E\le t_c\le3E$, and minimality of $c$
gives $t_{c-1}\le E$, which verifies the cutoff upper constraints. Thus
this point is feasible in the second family. The maximum of the two program
families is therefore $F$.

We next eliminate the two-large-modes family. At $B=t_b$, both selected
modes have allocation at most $B-E$; for mode $2$ this follows by
monotonicity of its deficit after $t_a$. All remaining modes have total
mass at most $T-2E$, giving $B\le2(B-E)+T-2E$ and
$E\le(T+B)/4$. The cutoff and $m_2\ge E$ also give
$E\le(T-Q)/2$. If $B\le T/3$ or $Q\ge T/3$, one of these bounds is at most
$T/3$. Thus only $b=c$ can give more than $T/3$, and the entire family
is bounded by $H$.
For attainment let $C=t_c$, $E=H$, and assign two identical modes total
mass $E$ each. At $C$, give each of them
\begin{equation}\label{eq:grid-H-witness}
 x=\max\{0,(C-T+2E)/2\};
\end{equation}
each of the other $n-2$ modes receives cumulative mass
$\min(C,T-2E)/(n-2)$ and terminal mass $(T-2E)/(n-2)$.
Interpolate on $[0,C]$ and $[C,T]$.
The inequalities $E\ge T/3$, $E\le C$, $4E\le T+C$, and
$T-C\le2E\le T-t_{c-1}$ verify every constraint with $a=b=c$.
This proves that the second family contributes exactly $H$.

For the remaining contribution only one distinguished mode is needed.
Return to the coverage argument with $m_2\le E<F$ and average all
modes other than mode $1$. Each averaged initial deficit still exceeds
$E$ at $t_a$. Their common terminal mass is at most the old $m_2$, so
the first eligible time for $1\to p$ moves weakly later. The initial
deficit of mode $1$ is nondecreasing and remains above $E$ there.
Thus this symmetrization preserves the obstruction, with a possibly new
$b\ge a$. Limits again cover boundary values; if the overall maximum
comes only from the other family, $H$ already covers it.

Write $M=m_1$, $x=A_1(A)$, $y=A_1(B)$. Every other mode has values
$(A-x)/(n-1)$, $(B-y)/(n-1)$, $(T-M)/(n-1)$. The failure constraints
are $x\ge(n-1)E-(n-2)A$ and $y\le B-E$. Monotonicity is equivalent to
\[
 0\le x\le y\le M,\qquad 0\le A-x\le B-y\le T-M.
\]
Together with the cutoff constraints and $M\ge T/n$, such states exist
if and only if $E\le A$, $nE\le(n-2)A+B$, and
\begin{equation}\label{eq:grid-M-interval}
\begin{split}
 \max\{T/n,\ T-A-E,\ (n-1)(E+Q)-(n-2)T,\ (n-1)E-(n-2)A\}
 \\
 \le M\le\min\{T-P-E,\ (n-1)(E+B)-(n-2)T\}.
\end{split}
\end{equation}
Necessity follows from the displayed monotonicity and cutoff inequalities.
For sufficiency choose any $M$ in this interval and set
\begin{equation}\label{eq:grid-xy-reconstruct}
 x=\max\{0,(n-1)E-(n-2)A,A-T+M\},\qquad
 y=\max\{x,B-T+M\}.
\end{equation}
The assumptions give $x\le\min(A,M,B-E)$ and
$y\le\min(M,B-E)$. They also give $y-x\le B-A$ and
$B-y\le T-M$, proving the required monotonicity. For example,
$B-T+M\le B-E$ follows from $M\le T-P-E\le T-E$;
$x\le B-E$ follows from $nE\le(n-2)A+B$ and this same upper bound.
These observations also cover $a=b$ and $b=N$.

Compare the four lower endpoints of \eqref{eq:grid-M-interval} with its
two upper endpoints. The comparisons involving $T-A-E$ versus
$T-P-E$, and $(n-1)(E+Q)-(n-2)T$ versus
$(n-1)(E+B)-(n-2)T$, hold automatically.
The comparison of $(n-1)E-(n-2)A$ with the second upper endpoint is
exactly the first inequality of \eqref{eq:grid-admissible}.
The remaining five comparisons, together with $E\ge T/3$, $E\le A$
and $nE\le(n-2)A+B$, are exactly $L_{ab}\le E\le U_{ab}$.
Maximization proves \eqref{eq:grid-minimax-formula};
\eqref{eq:grid-M-interval}--\eqref{eq:grid-xy-reconstruct} reconstruct
its first-family extremizer with at most three phases.
Each pair uses a fixed number of rational operations on the original grid
endpoints and $n$. The formulas have fixed algebraic depth, so intermediate
bit lengths are polynomial in the input length. Comparing $O(N^2)$
candidates proves the complexity claim.
\end{proof}

\subsection{Unit-grid consequences and the need for both families}
\begin{corollary}[Three modes, one switch]\label{cor:three-unit-one}
For $N\ge2$ unit cells,
\begin{equation}\label{eq:three-unit-one}
 F^{\mathrm{unit}}_{3,1}(N)=
 \begin{cases}
 k,&N=3k,\\
 k+\tfrac12,&N=3k+1\quad(k\ge1),\\
 k+\tfrac34,&N=3k+2\quad(k\ge0).
 \end{cases}
\end{equation}
For $N=1$ the value is $2/3$.
\end{corollary}
\begin{proof}
The one-cell case is Proposition~\ref{prop:no-switch}.
For $N=3k$, choosing the second-largest mass first until $k$ and the
largest mass last gives all three error terms at most $k$; three equal
masses give the matching omitted-mode lower bound.
For the other cases put $L=k+1$ and let $E$ be the claimed value. Direct
substitution gives
\begin{gather}\label{eq:residue-inequalities}
 E\ge N/3,\quad N-L\le2E,\quad N+L\le4E,\quad E+L\ge2N/3,\notag\\
 3E\ge2L,\qquad 3E+k+2L\ge2N.
\end{gather}
There are at most two masses greater than $E$. If there are two, their
allocations at $L$ sum to more than $L-N+2E\ge2(L-E)$.
One of them, used first, therefore has initial deficit at most $E$ at
$L$; the other, used last, has final deficit less than $N-E-L\le E$.
The omitted mass is at most $E$.
Otherwise let $q,r$ have the largest and second-largest masses. If
$m_q\ge N-E-k$, the order $r\to q$ at $k$ works. If not, then
$m_q<N-E-k$, while $m_q\ge N/3$ ensures that any $p\to q$ at $L$
has final deficit at most $E$. Such a schedule works if some $p\ne q$
has $A_p(L)\ge L-E$. In the remaining case
$A_q(L)>2E-L\ge L-E$, and
\[
 m_r\ge(N-m_q)/2>(E+k)/2\ge N-E-L.
\]
Thus $q\to r$ at $L$ works. These schedules omit only masses at most
$E$, establishing the upper bound.

For $N=3k+1$, use $k$ pure cells of mode $3$, one cell with rates
$(1/2,1/2,0)$, then $k$ pure cells of mode $1$ and $k$ of mode $2$.
Both selected large masses equal $E=k+1/2$, with allocation $1/2=L-E$
at $L$. Omitting either gives error $E$; selecting both gives final
deficit at least $E$ for switches at $\tau\le k$, and initial deficit
at least $E$ for $\tau\ge L$.
For $N=3k+2$, replace the mixed cell by $(1/4,1/4,1/2)$ and append
one final cell $(1/2,1/2,0)$. The two large masses are now $E=k+3/4$,
and each has allocation $1/4=L-E$ at $L$. The same argument applies:
early switches have final deficit at least $k+5/4>E$, and late switches
have initial deficit at least $E$. Empty pure blocks when $k=0$ cause
no change in the argument.
\end{proof}
The corrections to $N/3$ are respectively $0$, $1/6$, and $1/12$;
they are not a fixed half-cell correction.

\begin{example}[Five modes and nine unit cells]\label{ex:five-nine}
The exact value is $F^{\mathrm{unit}}_{5,1}(9)=17/5$. One extremizer has
rates
\[
\begin{array}{c|cc}
\text{segment}&\text{mode }1&\text{each of modes }2,3,4,5\\\hline
{[0,4]}&2/5&3/20\\
{[4,5]}&0&1/4\\
{[5,9]}&1/4&3/16
\end{array}
\]
Its terminal masses are $(13/5,8/5,8/5,8/5,8/5)$.
Every $p\to1$ with $p\ne1$ has final deficit at least $17/5$ when
$\tau\le3$, and initial deficit at least $4-3/5=17/5$ when $\tau\ge4$.
For $1\to2$, the analogous cutoff is between $4$ and $5$:
the final deficit at $4$ is $17/5$, and the initial deficit at $5$ is
$5-8/5=17/5$. Dominance proves the lower bound.
The upper bound also has a short direct proof. Sort the masses. If
$m_1\ge13/5$, use $2\to1$ at time $3$: the initial deficit is at most
$3$, the omitted mass is at most $m_3\le3$, and the final deficit is
$6-m_1\le17/5$. Otherwise $m_2\ge(9-m_1)/4>8/5$.
At time $4$ choose $p$ with $A_p(4)\ge4/5$ and finish in a largest-mass
mode $q\ne p$. The initial deficit is at most $16/5$, the final deficit
is $5-m_q<17/5$, and omitted masses are below $13/5$.
The value $17/5$ exceeds both the uniform-input value $16/5$ and the value
$3$ attainable by inputs supported on three modes. Hence neither family
attains the five-mode worst case.
\end{example}

The admissibility conditions in Theorem~\ref{thm:finite-one} also matter
on nonuniform grids. For
\[
 n=9,\qquad
 \mathcal T=(0,19/2,61/6,265/24,481/24,2669/120),
\]
a simplified first-family formula that retains only the upper bounds
$((n-2)A+B)/n$ and $((n-1)T-P-(n-1)Q)/n$ gives $2077/216$,
whereas the exact first-family maximum is $2593/270$.
Hence the other cutoff and feasibility conditions cannot be dropped.
The complete formula is exact and needs no LP solver. The supplementary
rational checker verifies these values, reconstructs extremizers, and
checks their absolute discrepancies by exhaustive schedule enumeration.
LP and MILP comparisons are supplementary numerical checks and are not
used in the proof.
